\documentclass{article}
\usepackage{PRIMEarxiv} % Core package for the PrimeAI Template
\usepackage{amsmath, amssymb, amsthm, bm, dcolumn} % Add amsthm here for the proof environment
\usepackage[numbers,sort&compress]{natbib} % Natbib for citations
\usepackage{graphicx} % For high-quality images
\usepackage{multicol} % Multi-column figures/tables
\usepackage{hyperref} % Hyperlinks
\usepackage{listings} % Code listings
\usepackage{authblk} % For structured affiliations
% Define theorem style (optional)
\theoremstyle{plain}
\newtheorem{theorem}{Theorem}
\renewcommand{\qedsymbol}{}

% Custom commands for primes and qubit encoding
\newcommand{\primeQubit}[1]{|p_{#1}\rangle}
\newcommand{\primeGate}[1]{U_{p_{#1}}}

\usepackage{wrapfig}
\usepackage[pscoord]{eso-pic}
\usepackage[fulladjust]{marginnote}
\reversemarginpar

% Typesetting improvements without footnote patching
\usepackage[protrusion=true, expansion=true, tracking=false]{microtype}
\microtypecontext{spacing=nonfrench}

% Line numbers
\usepackage[right]{lineno}

% Text layout - adjust as needed
\raggedright
\setlength{\parindent}{0.5cm}
\textwidth 5.25in 
\textheight 8.75in

% Set double spacing
\usepackage{setspace} 
\doublespacing

% Adjust width for specific content
\usepackage{changepage}

% Adjust caption style
\usepackage[aboveskip=1pt,labelfont=bf,labelsep=period,singlelinecheck=off]{caption}

% Remove brackets from references
\makeatletter
\renewcommand{\@biblabel}[1]{\quad#1.}
\makeatother

% Header, footer, and page numbers
\usepackage{lastpage,fancyhdr}
\pagestyle{fancy}
\fancyhf{}
\fancyhead[L]{Preprint - PrimeAI Enhanced Template}
\fancyfoot[C]{\scriptsize Multiplicity Theory © 2024 Ryan Van Gelder - Citizen Gardens \\ Licensed Under MIT and CC BY-NC-SA 4.0.}
\fancyfoot[R]{Page \thepage\ of \pageref{LastPage}}
\renewcommand{\footrule}{\hrule height 2pt \vspace{2mm}}

\title{Integrating Priyasha Mathur's Contributions with Multiplicity Theory}
\author{Ryan O. Van Gelder}
\affil{Citizen Gardens - The Foundation of Multiplicity \\ \texttt{info@citizengardens.org}}
\date{\today}

\begin{document}

\maketitle

\begin{abstract}
Priyasha Mathur’s contributions to prime encodings, dynamic systems, and algorithmic theory are integrated into Multiplicity Theory. Her research emphasizes the use of prime-based encoding for recursive systems, distributed networks, and computational frameworks, enhancing the adaptability and computational depth of Multiplicity Theory. This paper provides a structured mathematical framework for these integrations.
\end{abstract}

\section{Introduction}

Priyasha Mathur’s pioneering work in prime encodings, dynamic systems, and algorithmic theory offers significant advancements for Multiplicity Theory. This paper synthesizes her contributions into the principles of Multiplicity, focusing on prime-based encoding, recursive feedback, and their applications to distributed systems, cryptography, and quantum-inspired models.

\section{Prime Encodings in Multiplicity Theory}

\subsection{State Representation}
Prime-based encoding offers a unique mechanism to represent system states:
\begin{equation}
    S = \prod_{i=1}^n p_i^{x_i}, \quad x_i \in \mathbb{Z}^+,
\end{equation}
where \(p_i\) are primes encoding properties of the system, and \(x_i\) represents their respective multiplicities.

\subsection{Encoding Relationships}
Relationships between system components can be encoded using modular arithmetic:
\begin{equation}
    R(S_1, S_2) = S_1 \mod S_2,
\end{equation}
providing a compact representation of interactions.

\section{Dynamic Systems and Recursive Feedback}

\subsection{Recursive Feedback Loops}
Dynamic systems in Multiplicity Theory evolve under recursive rules with prime-based feedback:
\begin{equation}
    S_{t+1} = S_t + f(S_t, p_i),
\end{equation}
where \(f(S_t, p_i)\) introduces adjustments based on current states and prime encodings.

\subsection{Temporal Dynamics}
Time-evolution of the system can be expressed as:
\begin{equation}
    \frac{dS(t)}{dt} = \sum_{i=1}^n \phi(p_i, t),
\end{equation}
where \(\phi(p_i, t)\) captures the influence of \(p_i\) over time.

\section{Algorithmic Enhancements in Multiplicity}

\subsection{Prime-Based Sorting Algorithms}
Sorting algorithms leverage unique properties of primes:
\begin{equation}
    T(p_i, p_j) = 
    \begin{cases} 
    1 & \text{if } p_i < p_j, \\
    0 & \text{otherwise.}
    \end{cases}
\end{equation}

\subsection{Optimization via Prime Multiplicities}
Objective functions are optimized through prime encoding:
\begin{equation}
    \min_{x \in \mathbb{R}^n} f(x) = \sum_{i=1}^n p_i^{g_i(x)},
\end{equation}
where \(g_i(x)\) captures contributions to the system state.

\section{Distributed Networks}

\subsection{Prime-Based Graph Representation}
A graph \(G\) with vertices \(V\) and edges \(E\) can be encoded as:
\begin{equation}
    G = \bigcup_{(i,j) \in E} \phi(p_i, p_j),
\end{equation}
where \(\phi(p_i, p_j)\) encodes edge properties.

\subsection{Consensus Algorithms}
Consensus in distributed systems is achieved through modular arithmetic:
\begin{equation}
    S_\text{consensus} = \prod_{i=1}^n S_i \mod p_k.
\end{equation}

\section{Quantum-Inspired Systems}

\subsection{Prime-Encoded Quantum States}
Quantum states are represented using prime encodings:
\begin{equation}
    \psi(t) = \sum_{i=1}^n \alpha_i p_i^{\beta_i},
\end{equation}
where \(\alpha_i\) are amplitudes, and \(\beta_i\) encodes quantum properties.

\subsection{Tensor Networks for Multi-Dimensional Interactions}
Tensor networks model interactions between prime-based systems:
\begin{equation}
    T = \bigotimes_{i=1}^n p_i^{a_i}.
\end{equation}

\section{Unified Framework}
The integration of Priyasha Mathur’s contributions into Multiplicity Theory leads to the unified equation:
\begin{equation}
    E(t) = \int_X \left[ \sum_{i=1}^n p_i^{f_i(x)} + \phi(p_i, x, t) + \mathcal{F}(t) \right] dx,
\end{equation}
where:
\begin{itemize}
    \item \(p_i^{f_i(x)}\): Encodes system states and interactions.
    \item \(\phi(p_i, x, t)\): Represents temporal dynamics.
    \item \(\mathcal{F}(t)\): Captures recursive feedback mechanisms.
\end{itemize}

\section{Applications}

\subsection{Cryptography}
Prime encodings enhance cryptographic protocols by leveraging modular arithmetic for secure communications.

\subsection{Network Optimization}
Distributed networks are optimized using prime-based consensus algorithms and feedback mechanisms.

\subsection{Quantum Simulations}
Prime-encoded quantum states allow efficient modeling of entangled systems and recursive quantum operations.

\section{Conclusion}

Priyasha Mathur’s work on prime encodings, recursive feedback, and distributed systems greatly enhances Multiplicity Theory. These contributions enable robust applications in cryptography, network optimization, and quantum simulations, furthering the scope and adaptability of Multiplicity Theory.
\nocite{*}
\bibliographystyle{plain}
\bibliography{references}

\end{document}
